\chapter*{Preface}
\addcontentsline{toc}{chapter}{Preface}

Book 6 gave the plenum a connection, a metric, and a Ricci tensor, all built from entropy gradients, but left the resulting geometry ethically neutral: curvature there measured departure from flatness without asking what a system bent that way might be said to want. The present book removes that neutrality. Every metric developed here encodes an implicit valuation, and curvature increases specifically where a system cares most to persist. The universe, on the account this book develops, is organized not by blind dynamics smoothing indifferently but by selective smoothing, a continuous optimization of persistence under the fixed-volume constraint Book 6 already established.

\chapter*{Diagrammatic Overview}
\addcontentsline{toc}{chapter}{Diagrammatic Overview}

The passage from Book 6 to the present volume is a passage from thermodynamic connection to valuated geometry: the same entropic connection $\nabla^{(\SEntropy)}$ that Book 6 built from $\SEntropy$'s Hessian is retained here, but is now read alongside a preference potential $P_{\text{steady}}$ whose own Hessian gives the metric its ethical content. Nothing in the underlying connection changes; what changes is the interpretive weight this book asks the same formal structure to carry.

\chapter*{Reading Note}
\addcontentsline{toc}{chapter}{Reading Note}

Replace force with concern throughout this book: wherever a gradient appears, read it not as a cause pushing a system along but as an index of what that system is concerned to preserve. The metric of this book is an index of care, not a background structure the reader can set aside once its formal role is understood.

\part{Teleodynamics as Geometric Law}

\chapter{The Principle of Persistent Preference}

The teleodynamic principle governing this book is
\[
\frac{d}{dt} \grad_i \SEntropy = -\partial_i \PhiField + \alpha \, \grad_i \log P_{\text{steady}},
\]
tying the rate of change of entropy's gradient to the potential's own gradient, corrected by a term proportional to the log-gradient of a preference potential $P_{\text{steady}}$, weighted by a coupling constant $\alpha$. This preference potential is introduced here as the steady-state distribution a system's dynamics converge to when undisturbed, and persistence, on the reading this chapter proposes, is entropy-weighted optimization: a system persists not by resisting entropy production outright but by shaping where that production occurs so as to track $P_{\text{steady}}$'s own structure. The preference vector $\grad_i \log P_{\text{steady}}$ functions diagrammatically as a curvature bias, tilting the entropic connection of Book 6 in the direction the system is disposed to favor.

